\section{压力亏损怎样选出开放排气：从能量到重联率}
\label{sec:pressure-rate}
本节重建Liu等的低$\beta$、反平行、二维准稳态模型\citep{F589}。它提出一个比“Hall使速度变快”更具体的问题：为什么X点附近不能建立足以维持长薄层的压强？坐标取$x$为排出、$z$为入流、$y$为电流方向。本节$B_{xi}$、$B_{xe}$分别是在离子、电子扩散区上游边界的重联分量；$B_{x0}$为远上游场。三者不能相互替换。

\subsection{总Hall功为零，不妨碍离子获得能量}
在离子扩散区内、电子扩散区外，先忽略电子惯性和电子非理想压力项，得到
\begin{equation}
 \E=-\vel_e\times\B=-\vel_i\times\B+{\J\times\B\over en}.
\end{equation}
因此
\begin{align}
 \J\cdot\E&=\vel_i\cdot(\J\times\B)
             =en\vel_i\cdot\E_H,\\
 -en\vel_e\cdot\E&=0,
 \qquad \E_H={\J\times\B\over en}.
\end{align}
这里$\J\cdot\E_H=0$和$en\vel_i\cdot\E_H\ne0$并不矛盾：前者是两种带电粒子对这一电场的\emph{总}功，后者是其中一种的功。也不能由前者直接推出“无加热”：流体动能、压力张量和热流还会继续交换能量。

沿入流对称线，有$J_y\simeq\partial_zB_x/\mu_0$，Hall法向场为$E_{H,z}\simeq-J_yB_x/(en)$。取一条很窄的离子流管，其宽$\ell(z)$，沿$y$的长度为$\ell_y$。近不可压条件给
\begin{equation}
 |u_{iz}(z)|\ell(z)=C.
\end{equation}
从$d_i$到$d_e$，Hall法向场对离子的功率近似为
\begin{equation}
 W_i\simeq C\ell_y\int_{d_e}^{d_i}\partial_z{B_x^2\over2\mu_0}\,\dd z
 =C\ell_y{B_{xi}^2-B_{xe}^2\over2\mu_0}.
\end{equation}
这一步忽略$\partial_xB_z$对$J_y$的贡献，且仅统计建立入流方向压力的那部分能量。迎面进入的两束离子把有向运动变为以总体平均速度定义的$P_{izz}$。若压力增量主要在$zz$分量，压力能流
\begin{equation}
 \vct H_i={1\over2}\operatorname{tr}(\mathsf P_i)\vel_i+\mathsf P_i\cdot\vel_i
\end{equation}
的法向增量就是$\Delta H_{iz}\simeq(3/2)\Delta P_{izz}u_{iz}$。令能流增加等于上述电场功，消去$C\ell_y$，得
\begin{equation}
 \boxed{\Delta P_{izz}\simeq{2\over3}{B_{xi}^2-B_{xe}^2\over2\mu_0}.}
 \label{eq:remaining-pressure}
\end{equation}
系数$2/3$来自各向异性压力能流，并不是任意拟合系数；但“增加的能量主要进入$P_{izz}$”本身是此模型的物理近似。若粒子已显著热化、压缩或强热流不能忽略，就要重算能量预算。

\subsection{用相似张角连接两个尺度}
设离子和电子扩散区的半长分别为$L_i,L_e$，半厚近似$d_i,d_e$。假设分离线近似直线，则二者张角斜率相近：
\begin{equation}
 s\simeq{d_i\over L_i}\simeq{d_e\over L_e}.
\end{equation}
不可压连续性给$u_{{\rm in},s}/u_{{\rm out},s}\simeq s$；再假定两尺度的排出速度分别近于当地物种Alfvén速度。稳态Faraday定律要求$E_y$近似相同，因此
\begin{equation}
 s B_{xi}{B_{xi}\over\sqrt{\mu_0nm_i}}
 \simeq E_y\simeq
 s B_{xe}{B_{xe}\over\sqrt{\mu_0nm_e}},
 \qquad {B_{xe}^2\over B_{xi}^2}\simeq\sqrt{m_e\over m_i}.
 \label{eq:remaining-field-ratio}
\end{equation}
这个结果不是“电子比离子轻，所以直接用远上游电子Alfvén速度”——电子真正接收到的重联磁场也已降低。

接着用法向力平衡
\begin{equation}
 \partial_z\left({B^2\over2\mu_0}\right)+(\grad\cdot\mathsf P)_z
 ={(\B\cdot\grad)B_z\over\mu_0}.
\end{equation}
在两个尺度之间作有限差分，以平均场$(B_{xi}+B_{xe})/2$和同一张角估算右侧张力，得到
\begin{equation}
 s^2\simeq{B_{xi}^2-B_{xe}^2-2\mu_0\Delta P_{zz}\over(B_{xi}+B_{xe})^2}.
\end{equation}
代入式\eqref{eq:remaining-pressure}，并定义$r=B_{xe}/B_{xi}$，得
\begin{equation}
 \boxed{s^2\simeq{1\over3}{1-r\over1+r},\qquad
 r\simeq(m_e/m_i)^{1/4}.}
\end{equation}
压力只能补偿磁压下降的约三分之二，余下部分须由弯曲磁力线的张力补偿；这是开放排气的具体来源。

\subsection{上游场减弱和排出速度共同限制速率}
在离子扩散区外侧再作类似的磁压与张力平衡，令远近场差与弯曲项满足
$B_{x0}-B_{xi}\simeq s^2(B_{x0}+B_{xi})$，便有
\begin{equation}
 b\equiv{B_{xi}\over B_{x0}}\simeq{1-s^2\over1+s^2}.
\end{equation}
张角很大时，虽然几何排气口变宽，入流边界的可用磁场却减弱。再用含重联分量磁压反作用的排出估计
$u_{\rm out}\simeq bV_{Ai0}\sqrt{1-s^2}$，连续性给$u_{\rm in}\simeq s u_{\rm out}$，最终
\begin{equation}
 \boxed{R\equiv{E_y\over B_{x0}V_{Ai0}}
 \simeq s\left({1-s^2\over1+s^2}\right)^2\sqrt{1-s^2}.}
 \label{eq:remaining-rate}
\end{equation}
质量比400时基本式给$R\simeq0.172$，1836时约$0.157$；原文PIC算例实际约$0.075$，作者还需引用排出热压修正才能接近这个数。故该理论给出量级与几何机制，不能宣称原始基本式已无参数精确预测所有速率\citep{F589}。

\assume{这里同时用了低上游$\beta$、反平行、单X线、准稳态、近不可压、两尺度相似张角及局地速度估计。系统尺寸改变会影响磁通供给和排气形状；由本节不能证明全局平均率与系统尺寸无关。}

\section{电子单独重联：没有离子喷流，为什么仍受离子压力限制}
\label{sec:remaining-erec}
Liu等的另一模型\citep{F741}考虑整个活动区小于充分耦合的离子排气尺度。设离子电流远小于电子电流，$\J\simeq-en\vel_e$；用Ampère定律把电子Lorentz力化为磁应力：
\begin{equation}
 nm_e\vel_e\cdot\grad\vel_e
 ={\B\cdot\grad\B\over\mu_0}-\grad{B^2\over2\mu_0}
 -en\E-\grad\cdot\mathsf P_e.
\end{equation}
只取惯性与磁张力，会得到电子Alfvén喷流；实际还存在减速电子的面内$E_x$。取离X点距离$L_0$处的\emph{输运平台速度}$U$，排气半宽$\delta_0$、斜率$s=\delta_0/L_0$。它不是紧贴电子扩散区的最大喷流速度。

\subsection{从Hall电场估计平台速度}
稳态电场环量为零，磁场在排气口被电子流带动。沿环路估计得到
\begin{equation}
 E_{x,H}\sim U\Delta B_y{\delta_0\over L_0},\qquad
 \Delta B_y\sim\mu_0enU\delta_0,
\end{equation}
故$enE_{x,H}\sim\mu_0e^2n^2U^2\delta_0^2/L_0$。此估计把Hall磁场的对流反作用写成与$U^2$成比例的减速项。

周期小系统中，离子虽然没有形成有序喷流，磁岛内仍会积累离子压力。作者用经验函数$f(L_x)$表示其抵消可用张力的比例，把有限差分动量平衡整理为
\begin{equation}
 nm_eU^2+2\mu_0e^2n^2\delta_0^2U^2
 \simeq{B_{x0}^2\over\mu_0}(1-s^2)(1-f).
\end{equation}
将$d_e^2=m_e/(\mu_0ne^2)$、$d_i^2=m_i/(\mu_0ne^2)$代入，逐项除以$nm_iV_{Ai0}^2=B_{x0}^2/\mu_0$，得
\begin{equation}
 \boxed{{U\over V_{Ai0}}\simeq{d_i\over\delta_0}
 \sqrt{{(1-s^2)(1-f)\over2+(d_e/\delta_0)^2}}.}
 \label{eq:remaining-erec-speed}
\end{equation}
在$d_e\ll\delta_0\ll d_i$时，$U$大致与$\delta_0^{-1}$成正比；极小时电子惯性阻止这个表达式无限增长。$f\to1$表示反压将排气压住。原文选$f=\operatorname{sech}[L_x/(1.28d_i)]$；函数形状和长度$1.28d_i$并非从第一性原理唯一导出，应作为待检验闭合输入。

\subsection{还需要几何方程，不能单凭速度指定重联率}
电子尺度与整个小系统的连续性和同一$E_y$给
\begin{equation}
 {U\over V_{Ai0}}\simeq\left({B_{xe}\over B_{x0}}\right)^2
 \sqrt{m_i\over m_e}.
\end{equation}
沿用前节压力亏损估计，把该小系统上游直接取为$B_{x0}$，有
\begin{equation}
 {B_{xe}\over B_{x0}}\simeq{1-3s^2\over1+3s^2},\qquad
 \delta_0\simeq{L_xs\over4},\qquad R\simeq{sU\over V_{Ai0}}.
\end{equation}
实际计算时，对给定$L_x/d_i,m_i/m_e,f$，在$0<s<1/\sqrt3$中求使两种$U$相等的根，再求$R$。$L_0\simeq L_x/4$也是几何近似，不能删掉后声称率只由质量比决定。

原文真实质量比的周期PIC域为$L_x/d_i=1.28$到$7.68$，得到全局离子归一率$O(1)$。去掉反压所得约$4.28$是作者对更开放系统的上限估计，不是这些周期模拟的实测值。局部电子归一率还满足
\begin{equation}
 R_e={E_y\over B_{xe}V_{Ae}}
 =R\left({B_{x0}\over B_{xe}}\right)^2\sqrt{m_e\over m_i}.
\end{equation}
因此$B_{xe}$的取点若差一倍，$R_e$就差四倍。原文模型$R_e\simeq0.4$--$0.5$与在$z=d_e$取样的PIC约$0.1$确有此差异，不能将两种定义混写为理论和模拟完全相符。

\section{Full-F回旋流体：全密度、极化与求解器不能混为一谈}
\label{sec:remaining-fullf}
本章先前的FLR推导解释了导心密度与粒子密度不同。GREENY及其后续研究\citep{E473,L314}进一步保留密度的有限振幅变化，但仍是强导引场、低$\beta$、等温、二维板状、无磁曲率的约化模型。Full-F中的“全”指分布/密度展开的选择，不能解读成完整动理学。

\subsection{为何出现对数密度与正则动量}
等温压力$p_s=T_sN_s$除以密度后给
\begin{equation}
 {1\over N_s}\grad_\parallel p_s=T_s\grad_\parallel\ln N_s.
\end{equation}
电子平行电场又含$-\partial_tA_\parallel$。把它移到惯性项同侧，时间导数自然作用于$m_sU_s+q_sA_\parallel$。为使后面的系数可核对，本节采用$\rho_s=\sqrt{T_em_i}/(eB_0)$、$c_s=\sqrt{T_e/m_i}$，所有带下标$\mathrm{ph}$的量表示有量纲物理量：
\begin{gather}
 \vct x={\vct x_{\rm ph}\over\rho_s},\quad t={c_st_{\rm ph}\over\rho_s},\quad
 N_s={N_{s,\rm ph}\over N_0},\quad U_s={U_{s,\rm ph}\over c_s},\\
 \phi={e\phi_{\rm ph}\over T_e},\quad A_\parallel={A_{\parallel,\rm ph}\over B_0\rho_s},
 \qquad\beta={\mu_0N_0T_e\over B_0^2}.
\end{gather}
这里$\beta$没有通常热压比定义中的因子2。如此$q_sA_{\parallel,\rm ph}/(m_sc_s)=A_\parallel/\mu_s$，可定义带符号参数
\begin{equation}
 \mu_s={m_s\over Z_sm_i},\quad\tau_s={T_s\over Z_sT_e},\quad
 K_s=U_s+{\Gamma_{1s}A_\parallel\over\mu_s}.
\end{equation}
于是$\mu_e=-m_e/m_i<0$，$\tau_e=-1$；质量本身没有负号，有符号的是带电归一参数。忽略电子FLR、取$Z_i=1$后，所用无量纲系统的一种括号表示为
\begin{align}
 \partial_t n_e&=-[\phi,n_e]+[A_\parallel,n_eu_e],\\
 \partial_t N_i&=-[\psi_i,N_i]+[\Gamma_1A_\parallel,N_iU_i],\\
 \partial_t K_e&=-[\phi,K_e]+[A_\parallel,u_e^2/2]
                   -\mu_e^{-1}[A_\parallel,\ln n_e],\\
 \partial_t K_i&=-[\psi_i,K_i]+[\Gamma_1A_\parallel,U_i^2/2]
                   +\tau_i[\Gamma_1A_\parallel,\ln N_i].
\end{align}
其中$[f,g]=\partial_xf\partial_yg-\partial_yf\partial_xg$，$\psi_i$是包含回旋平均及极化修正的离子有效电势。按上述归一，磁扰动引起的平行导数是$-[A_\parallel,\cdot]$，在二维中仍非零。原文先用含$\beta$的磁势归一，随后将$\beta$吸收进$A_\parallel$；不能把吸收前后的变量混用。

令$N=1+\widetilde N$、$|\widetilde N|\ll1$，则
\begin{equation}
 \ln N=\widetilde N-\widetilde N^2/2+\cdots,\qquad
 [A,NU]=[A,U]+[A,\widetilde N U].
\end{equation}
小振幅截断会丢掉后一项等较高阶乘积；冻结极化中的密度则是另一个近似。只把一个方程中的$N$换成常数，并不自动获得整套一致的$\delta F$模型。

\subsection{极化方程是一条变系数椭圆约束}
为看清结构，取常导引场、单价离子，记$G=\sqrt{\Gamma_0}$，并将量纲常数吸收进$N_i$与$\phi$。Full-F极化关系可写为
\begin{equation}
 \mathcal A[N_i]\phi=\Gamma_1N_i-n_e,\qquad
 \mathcal A[N_i]=-\grad_\perp\cdot G N_i G\grad_\perp .
\end{equation}
在周期或使边界项消失的条件下，$G$自伴且$N_i>0$，因此
\begin{equation}
 \int\phi\mathcal A[N_i]\phi\,\dd V
 =\int N_i|G\grad_\perp\phi|^2\,\dd V\ge0.
\end{equation}
还可从这一约束重建相容的有效电势。定义辅助泛函
\begin{equation}
 \mathcal L_\phi=\int\left[(\Gamma_1N_i-n_e)\phi
             -{1\over2}N_i|G\grad_\perp\phi|^2\right]\dd V.
\end{equation}
先对$\phi$变分、分部积分，$\delta\mathcal L_\phi/\delta\phi=0$便恢复极化方程；再对$N_i$变分，利用$\Gamma_1$自伴，得到
\begin{equation}
 \psi_i={\delta\mathcal L_\phi\over\delta N_i}
       =\Gamma_1\phi-{1\over2}|G\grad_\perp\phi|^2.
\end{equation}
这一式说明有效电势中的二次项与极化算子配套，不能独立更改系数。这里给出的是采用$G=\sqrt{\Gamma_0}$时的相容重建；原稿不同位置的$G/\Gamma_1$与系数写法存在差异，复现其实现仍须核对代码版本，不能将本节重建冒充逐行代码说明。
这既解释椭圆求解器为什么可用，也给出正定极化能的检验。OB近似是把其中$N_i$换成$N_0$；长波近似则是$G\to1$。两者作用在不同对象上，可以分别开启，不能当成同一个开关。Full-F在长波极限\emph{仍}保留$\grad\cdot(N_i\grad\phi)$，除非再作OB近似。

与之对应的等温相对自由能应使用
\begin{equation}
 F_n=\sum_s\int T_s\left[N_s\ln{N_s\over N_{s0}}-N_s+N_{s0}\right]\dd V.
\end{equation}
写$N_s=N_{s0}(1+\epsilon_s)$，利用对数展开逐项相消，有
\begin{equation}
 F_n=\sum_s\int T_sN_{s0}
 \left({\epsilon_s^2\over2}-{\epsilon_s^3\over6}+\cdots\right)\dd V.
\end{equation}
所以$\delta F$二次自由能是Full-F的局地近似；它不是在任意振幅下都正确的能量。完整总能应把磁能$\int|\grad_\perp A_\parallel|^2/(2\mu_0)\,\dd V$仅加\emph{一次}，不能因原稿把它写在物种求和括号内而重复计入。

\subsection{密度和FLR为何使Ampère反演更难}
代入$u_e=K_e-A_\parallel/\mu_e$、$U_i=K_i-\Gamma_1A_\parallel$，Ampère关系成为
\begin{equation}
 -{\nabla_\perp^2 A_\parallel\over\beta}
 =-n_eK_e+{n_eA_\parallel\over\mu_e}
   +\Gamma_1(N_iK_i)-\Gamma_1(N_i\Gamma_1A_\parallel).
\end{equation}
常密度时每个Fourier波数可独立相除；变密度时$\widehat{NA}(\vct k)$是卷积，不能将$N(\vct x)$随意放进一个逐波数分母。实际核读的GREENY v1中，直接FFT路线冻结密度，SOR路线保留变密度但省去相关离子gyroaverage；这两条路线各有近似。后续正式代码是否已改变，需要核对具体版本和输入，不能仅凭“Full-F代码”名称确认所有项都已实现\citep{E473}。

\section{稳定本征值为什么仍可能暂态放大}
\label{sec:remaining-nonnormal}
Full-F研究提出非正规线性算子可能参与爆发前阶段\citep{L314}。为避免把线性代数名词误当机制证明，先用一个可以手算的例子。令
\begin{equation}
 {\dd\vct q\over\dd t}=A\vct q,\qquad
 A=\begin{pmatrix}-\gamma&K\\0&-\gamma\end{pmatrix},\quad\gamma>0.
\end{equation}
两个本征值都是$-\gamma$，长期都衰减；但
\begin{equation}
 e^{At}=e^{-\gamma t}\begin{pmatrix}1&Kt\\0&1\end{pmatrix}.
\end{equation}
若初值$(0,1)^T$，则$\|\vct q(t)\|^2=e^{-2\gamma t}(1+K^2t^2)$。例如$\gamma=1,K=4,t=0.5$时，范数平方为$5/e\simeq1.84$，比初值大，尽管两个本征值都稳定。原因是第二个变量暂时把能量式中的幅度传给第一个，之后才一起衰减。

真实等离子体的物理扰动能通常不是任意变量的Euclidean范数，而是
\begin{equation}
 W={1\over2}\vct q^\dagger Q\vct q,\quad Q=Q^\dagger>0,
 \qquad {\dd W\over\dd t}={1\over2}\vct q^\dagger(A^\dagger Q+QA)\vct q.
\end{equation}
令$\vct y=Q^{1/2}\vct q$，最大瞬时能量增长率受
$Q^{-1/2}(A^\dagger Q+QA)Q^{-1/2}$最大本征值控制；有限时间最大增益为
\begin{equation}
 G(t)=\left\|Q^{1/2}e^{At}Q^{-1/2}\right\|_2^2.
\end{equation}
因此必须先给出$Q$、边界条件、初值扰动和演化时间。条件数很大或伪谱伸入右半平面只说明有敏感性/增长可能，不能直接指定实际重联增长量；线性暂态与非线性磁岛快速合并也不是同一个阶段。

\section{多离子、负尘埃和中性粒子：从电流定义重新推导}
\label{sec:remaining-multispecies}
\subsection{冷离子为何能修正原来的Hall电场}
对若干正离子群$\alpha$，准中性与电流定义为
\begin{equation}
 n_e=\sum_\alpha Z_\alpha n_\alpha,\qquad
 \J=e\sum_\alpha Z_\alpha n_\alpha\vel_\alpha-en_e\vel_e.
\end{equation}
解出$\vel_e$，代回电子动量方程，逐项得到
\begin{equation}
 \E={\J\times\B\over en_e}
 -\sum_\alpha{Z_\alpha n_\alpha\over n_e}\vel_\alpha\times\B
 -{\grad\cdot\mathsf P_e\over en_e}-{m_e\over e}D_e\vel_e
 +{\vct R_e\over en_e}.
\end{equation}
当热离子已退磁而冷离子仍以$\E\times\B$漂移时，冷离子对流项不可忽略，它会部分抵消用总电流估计的Hall项。把所有离子压成一个标量温度流体，容易丢掉这种群间不同的磁化尺度。MMS分离线观测给出这种具体反例\citep{L227}；但它不是对X点重联电场的直接测量。

负尘埃加入后，应把相应$Z_d<0$保留在上述求和内，此时$n_i=n_e+|Z_d|n_d$。若$d_e\ll d_i\ll d_d$，内层相对运动是电子与离子，外层相对运动则可能是有效正电的轻组分与缓慢负尘埃。这允许嵌套、极性相反的Hall电流回路\citep{L283}。原来的单离子Hall项并没有一个可随意翻转的“符号开关”；必须从实际载流组分重新求电流和动量。

\subsection{中性质量改变的是有效Hall尺度，不是标准定义本身}
最小带电--中性系统还需
\begin{align}
 \rho_iD_i\vel_i&\simeq\J\times\B-\grad p_i-\alpha_{in}(\vel_i-\vel_n),\\
 \rho_nD_n\vel_n&=-\grad p_n+\alpha_{in}(\vel_i-\vel_n),
 \qquad\alpha_{in}=\rho_i\nu_{in}=\rho_n\nu_{ni}.
\end{align}
若研究时间长于两种组分的动量交换时间，近似$\vel_i\simeq\vel_n$，相加便得总惯性$\rho_{\rm tot}=\rho_i+\rho_n$，排出速度约为
$V_{A,\rm tot}=B/\sqrt{\mu_0\rho_{\rm tot}}$。Hall对流速度量级$B/(\mu_0en_iL)$与此速度相等时对应
\begin{equation}
 L_H={B/(\mu_0en_i)\over V_{A,\rm tot}}
 =d_i\sqrt{\rho_{\rm tot}\over\rho_i}.
\end{equation}
标准$d_i=\sqrt{m_i/(\mu_0n_ie^2)}$仍按离子密度定义。变大的是在强耦合近似下的有效转变尺度。如果$L_H$已超过装置可用宽度，离子层未必能按无限空间理论扩大；MRX观测到速度下降而层宽不明显变大，正需要检查这个有限边界问题\citep{L190}。

\section{平均Hall发电机与局地零功，以及移动X点的诊断}
\subsection{局地正交不等于平均投影为零}
对RFP，令横线表示平均，波浪表示涨落；密度涨落较小时，平行平均电动势包含
\begin{equation}
 \mathcal E_\parallel\simeq
 -\overline{\widetilde{\vel}\times\widetilde\B}_\parallel
 +{\overline{\widetilde\J\times\widetilde\B}_\parallel\over e\bar n}
 \simeq-\overline{\widetilde{\vel}_e\times\widetilde\B}_\parallel.
\end{equation}
投影方向是\emph{平均}磁场；$\widetilde\J\times\widetilde\B$不必垂直于平均磁场。它还以相关Lorentz力出现在平均动量方程里。因此Hall发电机既可帮助也可抵消MHD发电机，并与流的改变相关\citep{F798}。从局地$\J\cdot\E_H=0$取平均，实际上得到
\begin{equation}
 \bar\J\cdot\bar\E_H=-\overline{\widetilde\J\cdot\widetilde\E_H},
\end{equation}
表示不同尺度之间可交换能量，而非新产生焦耳热。

\subsection{非理想区的X点速度应由磁通鞍点求出}
选$\B_\perp=\grad\psi\times\hat{\vct z}$，则$\partial_t\psi=-E_z$（可加与空间无关的规范项）。X点满足$\grad\psi(\vct x_X(t),t)=0$。对时间求全导数：
\begin{equation}
 0=\grad\partial_t\psi+H_\psi\vel_X,
 \qquad\boxed{\vel_X=H_\psi^{-1}\grad E_z\big|_X},
\end{equation}
其中$H_\psi$是二维Hessian矩阵。只有鞍点非退化时可直接反演；生成或消灭磁岛附近应改用轨迹匹配，不能对近奇异矩阵硬求逆。

存在导引场时$\E\times\B/B^2$可在磁X点有定义，但非理想区域没有定理保证它等于上述鞍点速度。新湍流预印本采用这一漂移作为X点参考系\citep{F655}，应保留待验证标记。另需区分$D_e=\J\cdot(\E+\vel_e\times\B)$、压力应变和熵产生；局部$D_e$正峰并非不可逆加热证明，电流片计数也并非重联事件计数。
